\section{Conclusions and Future Work}
We introduced \texttt{PredEx}, the largest dataset for legal judgment prediction and explanation in this context, marking a significant advancement over previous datasets. Our research explored instruction tuning on Large Language Models (LLMs), showing promise in improving prediction accuracy and explanatory depth.

Looking ahead, our focus will be on training Indian Legal domain-specific Large Language Models. This approach will ensure that the models are ingrained with domain-specific knowledge, crucial for tasks like legal judgment prediction with explanations. Furthermore, we plan to undertake Supervised Fine-Tuning (SFT) on various downstream tasks, including the judgment prediction with explanation task. Another key objective will be to incorporate contextual understanding into the models to mitigate issues like hallucinated responses, a common challenge with generative models.

The question remains as we advance in this field: How ready is the State-of-the-Art to aid in explainable judgment prediction? Our future efforts aim to answer this question by refining the capabilities of AI in legal applications, making a significant contribution to the evolving field of AI-assisted legal judgment prediction and explanation. The ultimate goal is to develop AI tools that can not only alleviate the backlog in the Indian judiciary but also deliver justice efficiently and transparently. To further enhance the accuracy and reliability of our system, we plan to implement a Reinforcement Learning from Human Feedback (RLHF) pipeline. This pipeline aims to refine the model's predictions and explanations based on human feedback, ensuring that the outputs align more closely with expert legal understanding and reasoning. The inclusion of RLHF represents a significant advancement in developing AI systems for legal judgment prediction and explanation, as it allows for continuous improvement and adaptation based on real-world feedback and interactions.
